% Options for packages loaded elsewhere
\PassOptionsToPackage{unicode}{hyperref}
\PassOptionsToPackage{hyphens}{url}
\documentclass[
  11pt,
  a4paper,
]{article}
\usepackage{xcolor}
\usepackage[top=2.0cm,bottom=2.0cm,left=2.0cm,right=2.0cm]{geometry}
\usepackage{amsmath,amssymb}
\setcounter{secnumdepth}{-\maxdimen} % remove section numbering
\usepackage{iftex}
\ifPDFTeX
  \usepackage[T1]{fontenc}
  \usepackage[utf8]{inputenc}
  \usepackage{textcomp} % provide euro and other symbols
\else % if luatex or xetex
  \usepackage{unicode-math} % this also loads fontspec
  \defaultfontfeatures{Scale=MatchLowercase}
  \defaultfontfeatures[\rmfamily]{Ligatures=TeX,Scale=1}
\fi
\usepackage{lmodern}
\ifPDFTeX\else
  % xetex/luatex font selection
  \setmainfont[]{Times New Roman}
  \setsansfont[]{Arial}
  \setmonofont[]{Menlo}
\fi
% Use upquote if available, for straight quotes in verbatim environments
\IfFileExists{upquote.sty}{\usepackage{upquote}}{}
\IfFileExists{microtype.sty}{% use microtype if available
  \usepackage[]{microtype}
  \UseMicrotypeSet[protrusion]{basicmath} % disable protrusion for tt fonts
}{}
\makeatletter
\@ifundefined{KOMAClassName}{% if non-KOMA class
  \IfFileExists{parskip.sty}{%
    \usepackage{parskip}
  }{% else
    \setlength{\parindent}{0pt}
    \setlength{\parskip}{6pt plus 2pt minus 1pt}}
}{% if KOMA class
  \KOMAoptions{parskip=half}}
\makeatother
\setlength{\emergencystretch}{3em} % prevent overfull lines
\providecommand{\tightlist}{%
  \setlength{\itemsep}{0pt}\setlength{\parskip}{0pt}}
% 1-column header for CystoDS section report
\PassOptionsToPackage{table,xcdraw}{xcolor}
\usepackage[section]{placeins}
\usepackage{caption}
\usepackage{fancyhdr}
\usepackage{titlesec}
\usepackage{enumitem}
\usepackage{seqsplit}
\usepackage{booktabs}
\usepackage{graphicx}
\usepackage{microtype}
\usepackage{tabularx}
\usepackage{adjustbox}
\usepackage{float}

\definecolor{CystoBlue}{HTML}{1F4E79}
\definecolor{CystoTeal}{HTML}{2A7F62}
\definecolor{CystoGray}{HTML}{4A5568}
\definecolor{CystoDark}{HTML}{1A202C}

\graphicspath{{../../Docs/}{Docs/}{paper_assets/}{../paper_assets/}}

\captionsetup{font=small,labelfont={bf,color=CystoBlue}}
\captionsetup[table]{position=top,skip=3pt}
\captionsetup[figure]{position=bottom,skip=4pt}

\renewcommand{\figurename}{Hình}
\renewcommand{\tablename}{Bảng}

\titleformat{\section}{\Large\bfseries\color{CystoBlue}}{\thesection.}{0.4em}{}
\titleformat{\subsection}{\large\bfseries\color{CystoTeal}}{\thesubsection.}{0.4em}{}
\titleformat{\subsubsection}{\normalsize\bfseries\color{CystoDark}}{\thesubsubsection.}{0.4em}{}

\setlength{\parindent}{1.0em}
\setlength{\parskip}{0.35em plus 0.05em minus 0.05em}
\setlist{nosep,leftmargin=1.5em,topsep=0.3em}

\pagestyle{fancy}
\fancyhf{}
\fancyhead[L]{\small\color{CystoGray}CystoDS -- Phương Pháp Nghiên Cứu và Mô Hình Đề Xuất (3S-HFT)}
\fancyhead[R]{\small\color{CystoGray}CystoDS 2026}
\fancyfoot[C]{\small\color{CystoGray}\thepage}
\setlength{\headheight}{14pt}
\renewcommand{\headrulewidth}{0.35pt}

\AtBeginDocument{%
  \hypersetup{
    colorlinks=true,
    linkcolor=CystoBlue,
    urlcolor=CystoTeal,
    citecolor=CystoBlue,
  }%
}
\usepackage{bookmark}
\IfFileExists{xurl.sty}{\usepackage{xurl}}{} % add URL line breaks if available
\urlstyle{same}
\hypersetup{
  pdftitle={CystoDS: Phương Pháp Phân Cấp Ba Giai Đoạn Tuần Tự (3S-HFT v3.1)},
  hidelinks,
  pdfcreator={LaTeX via pandoc}}

\title{CystoDS: Phương Pháp Phân Cấp Ba Giai Đoạn Tuần Tự (3S-HFT v3.1)}
\usepackage{etoolbox}
\makeatletter
\providecommand{\subtitle}[1]{% add subtitle to \maketitle
  \apptocmd{\@title}{\par {\large #1 \par}}{}{}
}
\makeatother
\subtitle{Kiến trúc Swin-Tiny, Lịch trình Curriculum Warmup cho
Hierarchy Loss và Cơ chế Hierarchical Marginalization}
\author{}
\date{Báo cáo nghiên cứu 2026}

\begin{document}
\maketitle

\section{CystoDS: Phương Pháp Phân Cấp Ba Giai Đoạn Tuần Tự (3S-HFT
v3.1)}\label{cystods-phux1b0ux1a1ng-phuxe1p-phuxe2n-cux1ea5p-ba-giai-ux111oux1ea1n-tuux1ea7n-tux1ef1-3s-hft-v3.1}

\subsection{Kiến trúc Swin-Tiny, Lịch trình Curriculum Warmup cho
Hierarchy Loss và Cơ chế Hierarchical
Marginalization}\label{kiux1ebfn-truxfac-swin-tiny-lux1ecbch-truxecnh-curriculum-warmup-cho-hierarchy-loss-vuxe0-cux1a1-chux1ebf-hierarchical-marginalization}

\subsection{4. Phương pháp Đề xuất: CystoHier
(Methodology)}\label{phux1b0ux1a1ng-phuxe1p-ux111ux1ec1-xuux1ea5t-cystohier-methodology}

\subsubsection{4.1. Kiến trúc Swin-Tiny Đa Nhiệm Phân Cấp (Shared
Late-Stage)}\label{kiux1ebfn-truxfac-swin-tiny-ux111a-nhiux1ec7m-phuxe2n-cux1ea5p-shared-late-stage}

\begin{itemize}
\tightlist
\item
  \textbf{Backbone:} Swin-Tiny tiền huấn luyện ImageNet (28,3M tham số),
  chia 4 stages với kích thước cửa sổ \(7 \times 7\).
\item
  \textbf{Cấu hình Shared Late-Stage:} Toàn bộ 3 heads (Binary, Coarse,
  Fine) cùng nhận đặc trưng vector 768 chiều từ Stage 4 sau lớp
  LayerNorm.
\end{itemize}

\begin{verbatim}
[Phase 1: Representation Learning]
  • Open 100% Backbone + 3 Heads
  • Loss: BCE + CE + 0.10*SupCon
  • Warmup: w_hrc(t) = 0.25*min(1, t/12)
               │
               ▼
[Phase 2: Coarse Alignment]
  • Freeze Backbone & Fine/Binary
  • Loss: Coarse PP-SBS (Parameter Isolation)
               │
               ▼
[Phase 3: Fine Alignment]
  • Freeze Backbone & Coarse/Binary
  • Loss: Fine PP-SBS + 0.25*L_cf
\end{verbatim}

\subsubsection{4.2. Công thức toán học các hàm mất
mát}\label{cuxf4ng-thux1ee9c-touxe1n-hux1ecdc-cuxe1c-huxe0m-mux1ea5t-muxe1t}

\begin{itemize}
\tightlist
\item
  \textbf{Hàm Supervised Contrastive Loss (\(L_{\text{supcon}}\)):}
  \[L_{\text{supcon}} = \sum_{i \in I} \frac{-1}{|P(i)|} \sum_{p \in P(i)} \log \frac{e^{z_i \cdot z_p / \tau}}{\sum_{a \in A(i)} e^{z_i \cdot z_a / \tau}}\]
\item
  \textbf{Ràng buộc phân cấp Coarse-Fine (\(L_{\text{hierarchy}}\)):}
  \[L_{\text{cf}} = D_{\text{KL}}\left(P_{\text{coarse}} \,\parallel\, P_{\text{from\_fine}}\right)\]
  \[P_{\text{from\_fine}}(C) = \sum_{f \in \text{Children}(C)} P_{\text{fine}}(f)\]
\item
  \textbf{Patient-Prior Smoothed Balanced Softmax Loss (PP-SBS):}
  \[L_{\text{PP-SBS}}(y, \hat{z}) = -\log \frac{\pi_y \cdot e^{\hat{z}_y}}{\sum_{j} \pi_j \cdot e^{\hat{z}_j}}\]
  \[\pi_j = (\text{patients}_j + \epsilon)^{0{,}5}\]
\end{itemize}

\subsubsection{4.3. Lịch trình Curriculum Warmup cho Hierarchy
Loss}\label{lux1ecbch-truxecnh-curriculum-warmup-cho-hierarchy-loss}

\begin{itemize}
\tightlist
\item
  \textbf{Quy luật biến thiên trọng số ở Phase 1:}
  \[w_{\text{hrc}}(t) = 0{,}25 \times \min\left(1{,}0,\, \frac{t}{12}\right)\]
\item
  \textbf{Cơ chế tác động:}

  \begin{itemize}
  \tightlist
  \item
    \(t \le 4\): \(w_{\text{hrc}} \approx 0\), cho phép Backbone tự do
    định hình không gian biểu diễn cơ bản không bị ràng buộc phả hệ
    cưỡng bức.
  \item
    \(t > 4 \rightarrow 12\): Tăng dần đều để nắn chỉnh tính nhất quán
    phả hệ khi các đặc trưng đã chín muồi.
  \end{itemize}
\end{itemize}

\subsubsection{4.4. Suy Luận Đa Tầng Kết Hợp \& Tính Nhất Quán Phả
Hệ}\label{suy-luux1eadn-ux111a-tux1ea7ng-kux1ebft-hux1ee3p-tuxednh-nhux1ea5t-quuxe1n-phux1ea3-hux1ec7}

\begin{itemize}
\tightlist
\item
  \textbf{Định nghĩa Tính nhất quán Coarse-Fine (Consistency):}
  \[\text{Consistency} = \frac{1}{N}\sum_{i=1}^N \mathbb{I}\left[\arg\max P_{\text{coarse}}^{(i)} = \text{Parent}\left(\arg\max P_{\text{fine}}^{(i)}\right)\right]\]
\item
  \textbf{Công thức suy luận kết hợp:}
  \[P_{\text{ens}}(C) = \lambda P_{\text{coarse}}(C) + (1-\lambda) P_{\text{from\_fine}}(C)\]
\item
  \textbf{Tham số tối ưu:} \(\lambda = 0{,}25\) (75\% thông tin trích
  xuất từ phân phối của Fine Head).
\end{itemize}

\begin{center}\rule{0.5\linewidth}{0.5pt}\end{center}

\end{document}
